\documentclass[aps,prd,10pt,onecolumn,superscriptaddress,nofootinbib]{revtex4-2}
\usepackage[utf8]{inputenc}
\usepackage{amsmath}
\usepackage{amssymb}
\usepackage{booktabs}
\usepackage{longtable}
\usepackage{hyperref}

% Auto-synced version from CHANGELOG.md via Scripts/sync_version.py
% Auto-generated by sync_version.py — do not edit manually.
% Source of truth: CHANGELOG.md
\newcommand{\itsmVersion}{11.4.1}


\begin{document}

\title{Supplementary Material: Comprehensive SPARC MCMC Parameter Ledger}
\author{Brendon Boyd}
\email{brendon.boyd@itsm-cosmology.org}
\affiliation{Independent Researcher, Burswood, Western Australia, Australia}
\date{Version \itsmVersion\ | Compiled \today}

\maketitle

\section{Introduction}
This supplementary document catalogues the comprehensive parameter estimation results for the entire 175-galaxy sample from the SPARC database \cite{sparc}. Performed via parallelized Markov Chain Monte Carlo (MCMC) walkers (32 walkers $\times$ 3,000 steps per galaxy) within the Integrated Toroidal-Syntropic Model (ITSM) framework, these estimations resolve the stellar disk mass-to-light ratio ($\Upsilon_{\text{disk}}$), bulge mass-to-light ratio ($\Upsilon_{\text{bulge}}$), and the local Hubble expansion rate ($H_0$) tethered geometrically to the yield acceleration $a_0 = c H_0 / 2\pi$.

The ledger is categorized by galaxy name, optimized components, reduced chi-square ($\chi^2_\nu$), degrees of freedom (DoF), and quality categorization. Outliers that clipped the lower boundary wall ($H_0 \to 50.0 \ \text{km s}^{-1} \text{Mpc}^{-1}$) are designated as ``Clipped'' or ``Standard'' based on line-of-sight dust obscuration and turbulent compact kinematics, as analyzed in the main text \cite{itsm_main}.

\section{Methodology: Hierarchical Bayesian Inference}
The parameter extraction was executed using a hierarchical Markov Chain Monte Carlo (MCMC) ensemble sampler. The Integrated Toroidal-Syntropic Model (ITSM) imposes a rigid geometric dependency between the fundamental kinematic scale $a_0$ and the local Hubble parameter $H_0$, defined exactly by:
\begin{equation}
    a_0 = \frac{c H_0}{2\pi}
\end{equation}
where $c$ is the speed of light. Consequently, the rotation curves are evaluated not as independent systems, but as nodes within a globally constrained Hubble flow. The kinematic profile of each galaxy is governed by the Plenum Shear Ansatz:
\begin{equation}
    g_{\text{tot}} = g_{\text{bar}} + \frac{2}{3} \sqrt{g_{\text{bar}} a_0}
\end{equation}
where the pre-factor $2/3$ derives strictly from the geometric projection of the transverse shear plane onto the bulk $T^3$ manifold.

\subsection{MCMC Configuration and Priors}
For each galaxy, we implement an affine-invariant ensemble sampler with 32 walkers initialized across a tight ball near the Maximum A Posteriori (MAP) estimate. Convergence is enforced over 3,000 steps per galaxy (with a 600-step burn-in discard phase). The log-posterior $\ln P(\theta | D)$ is defined by the product of the uniform priors and the Gaussian likelihood:
\begin{equation}
    \ln \mathcal{L} = -0.5 \sum_i \left( \frac{V_{\text{obs},i} - V_{\text{ITSM},i}}{\sigma_{V,i}} \right)^2
\end{equation}
To ensure the physics remains scale-invariant, broad uniform priors are assigned to the baryonic mass-to-light ratios ($\Upsilon_{\text{disk}} \in [0.01, 8.0]$, $\Upsilon_{\text{bulge}} \in [0.00, 8.0]$). Convergence is verified per galaxy via Gelman-Rubin diagnostics and integrated autocorrelation times ($\tau \ll N_{\text{steps}}$).

\section{Comprehensive SPARC Parameter Ledger}
\begin{longtable}{lcccccc}
\toprule
Galaxy & $\Upsilon_{\text{disk}}$ & $\Upsilon_{\text{bulge}}$ & $H_0$ & $\chi^2_\nu$ & DoF & Quality \\
\midrule
CamB & 0.013 & 3.942 & 64.1 & 1.63 & 6 & High-Fi \\
D512-2 & 0.278 & 3.754 & 71.9 & 2.27 & 1 & High-Fi \\
D564-8 & 0.259 & 3.949 & 71.4 & 1.20 & 3 & High-Fi \\
D631-7 & 0.114 & 4.237 & 71.6 & 1.63 & 13 & High-Fi \\
DDO064 & 0.010 & 3.954 & 52.3 & 7.29 & 11 & Standard \\
DDO154 & 0.349 & 4.027 & 72.1 & 4.70 & 9 & High-Fi \\
DDO161 & 0.040 & 3.941 & 69.7 & 2.39 & 28 & High-Fi \\
DDO168 & 0.140 & 4.040 & 71.4 & 1.52 & 7 & High-Fi \\
DDO170 & 2.643 & 4.006 & 73.2 & 2.28 & 5 & High-Fi \\
ESO079-G014 & 0.012 & 3.808 & 63.6 & 1.09 & 12 & High-Fi \\
ESO116-G012 & 0.012 & 3.934 & 65.4 & 1.19 & 12 & High-Fi \\
ESO444-G084 & 0.043 & 4.002 & 71.4 & 3.20 & 4 & High-Fi \\
ESO563-G021 & 0.011 & 4.096 & 65.5 & 0.81 & 27 & High-Fi \\
F561-1 & 0.164 & 3.797 & 71.2 & 0.95 & 3 & High-Fi \\
F563-1 & 0.070 & 4.113 & 70.7 & 1.08 & 14 & High-Fi \\
F563-V1 & 0.142 & 3.923 & 70.8 & 0.94 & 3 & High-Fi \\
F563-V2 & 0.011 & 3.801 & 57.9 & 3.40 & 7 & High-Fi \\
F565-V2 & 0.158 & 4.063 & 70.4 & 1.61 & 4 & High-Fi \\
F567-2 & 0.227 & 3.992 & 71.8 & 2.36 & 2 & High-Fi \\
F568-1 & 0.013 & 4.011 & 65.6 & 1.60 & 9 & High-Fi \\
F568-3 & 0.015 & 3.960 & 65.7 & 0.66 & 15 & High-Fi \\
F568-V1 & 0.012 & 3.961 & 59.3 & 1.71 & 12 & High-Fi \\
F571-8 & 0.010 & 4.118 & 62.2 & 3.82 & 10 & High-Fi \\
F571-V1 & 0.367 & 4.026 & 73.4 & 1.28 & 4 & High-Fi \\
F574-1 & 0.015 & 3.937 & 66.7 & 0.76 & 11 & High-Fi \\
F574-2 & 0.147 & 3.941 & 71.1 & 2.43 & 2 & High-Fi \\
F579-V1 & 0.012 & 3.910 & 63.2 & 1.21 & 11 & High-Fi \\
F583-1 & 0.011 & 4.006 & 55.1 & 2.05 & 22 & High-Fi \\
F583-4 & 0.011 & 4.106 & 57.3 & 2.39 & 9 & High-Fi \\
IC2574 & 0.056 & 3.948 & 71.7 & 0.40 & 31 & High-Fi \\
IC4202 & 0.013 & 0.000 & 53.6 & 4.62 & 29 & High-Fi \\
KK98-251 & 0.020 & 4.128 & 68.0 & 1.34 & 12 & High-Fi \\
NGC0024 & 0.011 & 3.983 & 55.7 & 1.37 & 26 & High-Fi \\
NGC0055 & 0.301 & 3.984 & 73.7 & 0.42 & 18 & High-Fi \\
NGC0100 & 0.011 & 3.975 & 58.4 & 1.35 & 18 & High-Fi \\
NGC0247 & 0.061 & 3.939 & 73.2 & 0.23 & 23 & High-Fi \\
NGC0289 & 0.022 & 4.031 & 71.7 & 0.53 & 25 & High-Fi \\
NGC0300 & 0.045 & 4.178 & 71.2 & 0.33 & 22 & High-Fi \\
NGC0801 & 0.013 & 4.032 & 73.4 & 0.84 & 10 & High-Fi \\
NGC0891 & 0.012 & 0.010 & 68.8 & 1.40 & 15 & High-Fi \\
NGC1003 & 0.066 & 3.887 & 72.8 & 0.56 & 33 & High-Fi \\
NGC1090 & 0.011 & 4.173 & 64.9 & 1.10 & 21 & High-Fi \\
NGC1705 & 0.013 & 4.114 & 68.3 & 1.02 & 11 & High-Fi \\
NGC2366 & 0.011 & 4.051 & 56.9 & 2.82 & 23 & High-Fi \\
NGC2403 & 0.010 & 4.077 & 57.1 & 1.14 & 70 & High-Fi \\
NGC2683 & 0.099 & 4.113 & 72.3 & 0.48 & 8 & High-Fi \\
NGC2841 & 0.026 & 4.006 & 63.0 & 0.65 & 47 & High-Fi \\
NGC2903 & 0.011 & 3.860 & 69.6 & 0.87 & 31 & High-Fi \\
NGC2915 & 0.021 & 3.887 & 70.0 & 0.72 & 27 & High-Fi \\
NGC2955 & 0.011 & 0.000 & 59.9 & 2.53 & 21 & High-Fi \\
NGC2976 & 0.010 & 3.956 & 52.8 & 3.14 & 24 & High-Fi \\
NGC2998 & 0.010 & 4.002 & 58.6 & 4.58 & 10 & High-Fi \\
NGC3109 & 0.018 & 4.106 & 68.4 & 1.27 & 22 & High-Fi \\
NGC3198 & 0.011 & 3.926 & 62.8 & 0.57 & 40 & High-Fi \\
NGC3521 & 0.011 & 4.056 & 65.1 & 0.83 & 38 & High-Fi \\
NGC3726 & 0.128 & 4.045 & 72.8 & 0.32 & 9 & High-Fi \\
NGC3741 & 0.047 & 4.076 & 71.1 & 4.66 & 18 & High-Fi \\
NGC3769 & 0.128 & 4.197 & 74.3 & 0.44 & 9 & High-Fi \\
NGC3877 & 0.018 & 3.945 & 72.7 & 0.42 & 10 & High-Fi \\
NGC3893 & 0.052 & 3.998 & 75.2 & 0.34 & 7 & High-Fi \\
NGC3917 & 0.022 & 3.933 & 72.3 & 0.33 & 14 & High-Fi \\
NGC3949 & 0.043 & 3.924 & 73.9 & 0.53 & 4 & High-Fi \\
NGC3953 & 0.104 & 3.784 & 72.9 & 0.40 & 5 & High-Fi \\
NGC3972 & 0.028 & 4.026 & 73.1 & 0.52 & 7 & High-Fi \\
NGC3992 & 0.931 & 3.745 & 74.9 & 0.34 & 6 & High-Fi \\
NGC4010 & 0.018 & 3.893 & 70.3 & 0.60 & 9 & High-Fi \\
NGC4013 & 0.286 & 3.945 & 65.1 & 0.35 & 33 & High-Fi \\
NGC4051 & 0.076 & 3.929 & 73.2 & 0.57 & 4 & High-Fi \\
NGC4068 & 0.015 & 3.984 & 66.2 & 3.13 & 3 & High-Fi \\
NGC4085 & 0.019 & 4.127 & 72.4 & 0.94 & 4 & High-Fi \\
NGC4088 & 0.029 & 4.050 & 73.1 & 0.37 & 9 & High-Fi \\
NGC4100 & 0.053 & 4.028 & 73.3 & 0.15 & 21 & High-Fi \\
NGC4138 & 0.202 & 4.067 & 71.5 & 1.32 & 4 & High-Fi \\
NGC4157 & 0.021 & 4.042 & 73.3 & 0.27 & 14 & High-Fi \\
NGC4183 & 0.034 & 4.201 & 70.8 & 0.39 & 20 & High-Fi \\
NGC4214 & 0.013 & 4.181 & 68.5 & 0.91 & 11 & High-Fi \\
NGC4217 & 0.027 & 0.000 & 51.7 & 14.75 & 16 & Standard \\
NGC4389 & 0.019 & 4.145 & 73.4 & 1.20 & 3 & High-Fi \\
NGC4559 & 0.017 & 4.092 & 70.4 & 0.41 & 29 & High-Fi \\
NGC5005 & 0.010 & 0.000 & 56.8 & 4.00 & 15 & High-Fi \\
NGC5033 & 0.012 & 0.001 & 64.6 & 1.50 & 19 & High-Fi \\
NGC5055 & 0.012 & 4.030 & 71.0 & 0.42 & 25 & High-Fi \\
NGC5371 & 0.027 & 4.060 & 72.4 & 0.21 & 16 & High-Fi \\
NGC5585 & 0.010 & 3.927 & 56.3 & 2.84 & 21 & High-Fi \\
NGC5907 & 0.173 & 3.848 & 73.7 & 0.22 & 16 & High-Fi \\
NGC5985 & 0.036 & 4.176 & 68.0 & 0.42 & 30 & High-Fi \\
NGC6015 & 0.011 & 4.167 & 66.9 & 0.56 & 41 & High-Fi \\
NGC6195 & 0.016 & 0.000 & 57.4 & 3.87 & 20 & High-Fi \\
NGC6503 & 0.015 & 3.955 & 71.2 & 0.33 & 28 & High-Fi \\
NGC6674 & 0.036 & 3.840 & 68.9 & 0.75 & 12 & High-Fi \\
NGC6789 & 0.011 & 4.080 & 60.1 & 17.13 & 1 & Standard \\
NGC6946 & 0.010 & 0.000 & 52.5 & 3.88 & 55 & High-Fi \\
NGC7331 & 0.021 & 4.040 & 71.4 & 0.16 & 33 & High-Fi \\
NGC7793 & 0.010 & 4.034 & 52.1 & 3.19 & 43 & High-Fi \\
NGC7814 & 0.016 & 0.000 & 52.9 & 11.75 & 15 & Standard \\
PGC51017 & 0.029 & 4.209 & 70.1 & 1.82 & 3 & High-Fi \\
UGC00128 & 0.051 & 4.123 & 70.7 & 0.95 & 19 & High-Fi \\
UGC00191 & 0.025 & 3.880 & 70.4 & 0.92 & 6 & High-Fi \\
UGC00634 & 3.788 & 3.853 & 74.6 & 3.28 & 1 & High-Fi \\
UGC00731 & 0.328 & 4.026 & 72.2 & 3.91 & 9 & High-Fi \\
UGC00891 & 1.916 & 3.916 & 72.3 & 2.66 & 2 & High-Fi \\
UGC01230 & 0.026 & 3.884 & 69.6 & 1.24 & 8 & High-Fi \\
UGC01281 & 0.010 & 4.024 & 51.1 & 9.55 & 22 & Standard \\
UGC02023 & 0.030 & 3.887 & 70.5 & 2.10 & 2 & High-Fi \\
UGC02259 & 0.112 & 4.083 & 72.5 & 0.77 & 5 & High-Fi \\
UGC02455 & 0.018 & 3.807 & 72.4 & 0.92 & 5 & High-Fi \\
UGC02487 & 0.277 & 3.910 & 69.7 & 0.65 & 14 & High-Fi \\
UGC02885 & 0.020 & 0.008 & 66.9 & 1.42 & 16 & High-Fi \\
UGC02916 & 0.011 & 0.000 & 50.3 & 24.83 & 40 & Clipped \\
UGC02953 & 0.010 & 0.000 & 50.2 & 16.82 & 112 & Clipped \\
UGC03205 & 0.010 & 0.000 & 50.9 & 7.32 & 45 & Clipped \\
UGC03546 & 0.011 & 0.000 & 52.7 & 5.58 & 27 & Standard \\
UGC03580 & 0.010 & 0.000 & 50.6 & 8.68 & 44 & Clipped \\
UGC04278 & 0.011 & 4.137 & 53.4 & 2.48 & 22 & High-Fi \\
UGC04305 & 0.012 & 4.010 & 59.6 & 1.35 & 19 & High-Fi \\
UGC04325 & 0.037 & 4.052 & 70.6 & 0.79 & 5 & High-Fi \\
UGC04483 & 0.012 & 3.817 & 60.0 & 4.75 & 5 & High-Fi \\
UGC04499 & 0.137 & 4.179 & 71.9 & 0.93 & 6 & High-Fi \\
UGC05005 & 0.028 & 4.019 & 69.2 & 1.29 & 8 & High-Fi \\
UGC05253 & 0.010 & 0.000 & 50.2 & 30.12 & 70 & Clipped \\
UGC05414 & 0.054 & 3.982 & 71.9 & 1.11 & 3 & High-Fi \\
UGC05716 & 0.288 & 3.923 & 71.3 & 2.31 & 9 & High-Fi \\
UGC05721 & 0.011 & 4.045 & 57.9 & 2.48 & 20 & High-Fi \\
UGC05750 & 0.012 & 3.818 & 60.2 & 2.07 & 8 & High-Fi \\
UGC05764 & 0.050 & 4.039 & 71.9 & 2.57 & 7 & High-Fi \\
UGC05829 & 0.025 & 3.999 & 67.6 & 1.96 & 8 & High-Fi \\
UGC05918 & 0.041 & 4.150 & 69.3 & 0.93 & 5 & High-Fi \\
UGC05986 & 0.023 & 3.918 & 70.1 & 0.50 & 12 & High-Fi \\
UGC05999 & 2.619 & 4.033 & 72.7 & 2.03 & 2 & High-Fi \\
UGC06399 & 0.040 & 4.020 & 70.6 & 0.68 & 6 & High-Fi \\
UGC06446 & 0.039 & 3.892 & 71.4 & 1.20 & 14 & High-Fi \\
UGC06614 & 0.177 & 0.011 & 53.6 & 11.17 & 10 & Standard \\
UGC06628 & 0.070 & 4.072 & 70.9 & 0.87 & 4 & High-Fi \\
UGC06667 & 0.049 & 4.009 & 71.8 & 1.33 & 6 & High-Fi \\
UGC06786 & 0.010 & 0.000 & 50.2 & 30.03 & 42 & Clipped \\
UGC06787 & 0.010 & 0.000 & 50.1 & 46.07 & 68 & Clipped \\
UGC06818 & 0.060 & 3.902 & 71.2 & 0.67 & 5 & High-Fi \\
UGC06917 & 0.199 & 4.057 & 73.4 & 0.37 & 8 & High-Fi \\
UGC06923 & 0.099 & 3.880 & 72.4 & 0.81 & 3 & High-Fi \\
UGC06930 & 0.171 & 4.023 & 73.1 & 0.62 & 7 & High-Fi \\
UGC06973 & 0.036 & 4.087 & 73.2 & 0.62 & 6 & High-Fi \\
UGC06983 & 0.196 & 3.851 & 74.0 & 0.58 & 14 & High-Fi \\
UGC07089 & 0.039 & 3.939 & 70.2 & 0.45 & 9 & High-Fi \\
UGC07125 & 0.222 & 3.994 & 72.6 & 1.10 & 10 & High-Fi \\
UGC07151 & 0.021 & 3.930 & 70.6 & 0.68 & 8 & High-Fi \\
UGC07232 & 0.016 & 3.890 & 67.8 & 5.66 & 1 & Standard \\
UGC07261 & 0.047 & 3.878 & 70.9 & 1.04 & 4 & High-Fi \\
UGC07323 & 0.019 & 3.829 & 69.6 & 0.75 & 7 & High-Fi \\
UGC07399 & 0.055 & 4.143 & 73.1 & 0.61 & 7 & High-Fi \\
UGC07524 & 0.022 & 3.849 & 69.0 & 0.46 & 28 & High-Fi \\
UGC07559 & 0.030 & 3.817 & 70.2 & 1.33 & 4 & High-Fi \\
UGC07577 & 0.013 & 4.070 & 60.8 & 2.98 & 6 & High-Fi \\
UGC07603 & 0.027 & 4.094 & 71.3 & 0.70 & 9 & High-Fi \\
UGC07608 & 0.064 & 4.053 & 72.0 & 1.31 & 5 & High-Fi \\
UGC07690 & 0.058 & 4.123 & 73.6 & 0.87 & 4 & High-Fi \\
UGC07866 & 0.023 & 3.997 & 69.2 & 1.31 & 4 & High-Fi \\
UGC08286 & 0.021 & 4.019 & 70.7 & 0.74 & 14 & High-Fi \\
UGC08490 & 0.018 & 3.962 & 71.8 & 0.78 & 27 & High-Fi \\
UGC08550 & 0.051 & 3.675 & 71.7 & 1.13 & 8 & High-Fi \\
UGC08699 & 0.010 & 0.000 & 51.2 & 6.90 & 38 & Standard \\
UGC08837 & 0.033 & 4.020 & 70.0 & 1.37 & 5 & High-Fi \\
UGC09037 & 0.104 & 4.054 & 73.4 & 0.23 & 19 & High-Fi \\
UGC09133 & 0.010 & 0.000 & 51.3 & 5.18 & 65 & Standard \\
UGC09992 & 0.130 & 4.055 & 71.6 & 1.42 & 2 & High-Fi \\
UGC10310 & 0.075 & 4.019 & 70.1 & 1.03 & 4 & High-Fi \\
UGC11455 & 0.011 & 3.752 & 66.8 & 0.55 & 33 & High-Fi \\
UGC11557 & 0.012 & 3.905 & 64.6 & 1.49 & 9 & High-Fi \\
UGC11820 & 0.012 & 4.067 & 61.0 & 3.98 & 7 & High-Fi \\
UGC11914 & 0.010 & 0.000 & 51.1 & 5.98 & 62 & Standard \\
UGC12506 & 0.013 & 3.861 & 67.1 & 0.48 & 28 & High-Fi \\
UGC12632 & 0.040 & 3.911 & 69.6 & 0.82 & 12 & High-Fi \\
UGC12732 & 0.089 & 3.898 & 72.0 & 1.36 & 13 & High-Fi \\
UGCA281 & 0.011 & 4.065 & 54.8 & 8.94 & 4 & Standard \\
UGCA442 & 0.086 & 3.956 & 71.6 & 3.07 & 5 & High-Fi \\
UGCA444 & 0.132 & 3.963 & 72.5 & 5.18 & 33 & Standard \\
\bottomrule
\end{longtable}


\bibliographystyle{apsrev4-2}
\bibliography{../references}

\end{document}
